\section{Defects}\label{app:E.10}

A carbon pair with a bound hydrogen, the T center of silicon, emits its zero-phonon line at \SI{935}{meV}; a hybrid-functional calculation in a 512-atom supercell puts it at \SI{985}{meV}~\cite{dhaliah2022first}. Displacing a silicon atom from its site takes between 17.5 and \SI{63.5}{eV}, depending on the direction in which it is struck~\cite{kadribasic2018directional}. Both numbers are defect properties that enter the book: the first in the color centers and dopants that serve as targets and readouts (\cref{sec:7.4,sec:16.6}), the second in the energy that a nuclear recoil stores in the lattice (\cref{chap:19}). The effective-mass theory of shallow donors, which no supercell reaches, is in \cref{sec:7.4,app:D.23}.

\subsection{Formation energies and charge states}

The formation energy of a defect $D$ in charge state $z_{\mathrm{D}}$, in units of $e$, compares two supercells, with and without the defect,
\begin{equation}
  E^{\mathrm{f}}_{\mathrm{D}}(z_{\mathrm{D}}) = E_{\mathrm{D}}(z_{\mathrm{D}}) - E_{\mathrm{host}} - \sum_s n_s\,\mu^{\mathrm{at}}_s + z_{\mathrm{D}}\bigl(\varepsilon_{\mathrm{VBM}} + \mu_{\mathrm{ch}}\bigr) + E_{\mathrm{corr}} ,
  \label{eq:E.10-form}
\end{equation}
where $E_{\mathrm{D}}$ and $E_{\mathrm{host}}$ are the total energies of the defective and the host supercell, $n_s$ atoms of species $s$ are added ($n_s > 0$) or removed ($n_s < 0$) and exchanged with a reservoir at chemical potential $\mu^{\mathrm{at}}_s$, the electron chemical potential $\mu_{\mathrm{ch}}$ is measured from the valence maximum $\varepsilon_{\mathrm{VBM}}$, and $E_{\mathrm{corr}}$ is the finite-size correction of the next subsection~\cite[Eq.~(3)]{kim2020quickstart}. The atomic chemical potentials are bounded by the stability of the host against its elements and against competing phases, and $\mu_{\mathrm{ch}}$ ranges over the gap~\cite[Sect.~3]{kim2020quickstart}. Where two charge states have equal formation energy, $E^{\mathrm{f}}_{\mathrm{D}}(z_{\mathrm{D}}) = E^{\mathrm{f}}_{\mathrm{D}}(z'_{\mathrm{D}})$, \cref{eq:E.10-form} gives the thermodynamic transition level
\begin{equation}
  \mu_{\mathrm{ch}}(z_{\mathrm{D}}/z'_{\mathrm{D}}) = \frac{E^{\mathrm{f}}_{\mathrm{D}}(z_{\mathrm{D}})\bigr|_{\mu_{\mathrm{ch}}=0} - E^{\mathrm{f}}_{\mathrm{D}}(z'_{\mathrm{D}})\bigr|_{\mu_{\mathrm{ch}}=0}}{z'_{\mathrm{D}} - z_{\mathrm{D}}} ,
  \label{eq:E.10-level}
\end{equation}
the electron chemical potential at which the defect changes charge with full relaxation of the lattice. It is the quantity that photoluminescence quenching and capture experiments measure, and it differs from a Kohn--Sham eigenvalue, which carries no relaxation and no correction.

The defect is relaxed in a supercell as close to cubic as possible, built from the conventional cell, at the volume of the host~\cite[Sect.~3]{kim2020quickstart}. The work lies in the number of charge and spin states, which must all be tried, since a defect in a nonmagnetic host may have a spin-polarized ground state. The gap problem of semilocal functionals is more serious for defects than for bulk bands, because defect levels sit inside the gap and their position relative to both band edges matters; hybrid functionals, which give silicon a gap of \SI{1.114}{eV}~\cite{dhaliah2022first}, are the standard for levels and optical energies.

\subsection{Charged defects in a periodic cell}

A charged defect in a supercell is a periodic array of charges, neutralized by a uniform background so that the electrostatic energy is finite. Both the interaction with the images and the background are artifacts. The leading term is the energy of a simple-cubic lattice of point charges $z_{\mathrm{D}}e$ in a compensating background and a dielectric medium. In Heaviside--Lorentz units, with $e^2 = 4\pi\alpha$, the Coulomb energy of two charges $e$ at distance $r$ in vacuum is $\alpha/r$, and the Madelung energy per charge of the array is $-z_{\mathrm{D}}^2\alpha\,\alpha_{\mathrm{M}}/(2\varepsilon L_{\mathrm{sc}})$, with $\alpha_{\mathrm{M}} = 2.8373$ the Madelung constant of the simple-cubic lattice~\cite[Sect.~II]{komsa2009comparison}. The correction removes it:
\begin{equation}
  E_{\mathrm{corr}} \simeq \frac{z_{\mathrm{D}}^2\,\alpha\,\alpha_{\mathrm{M}}}{2\varepsilon\,L_{\mathrm{sc}}} + O(L_{\mathrm{sc}}^{-3}) ,
  \label{eq:E.10-mp}
\end{equation}
where the next term, of order $z_{\mathrm{D}}L_{\mathrm{sc}}^{-3}$ times the quadrupole moment of the defect charge, comes from the shape of that charge; this is the Makov--Payne form, quoted in Gaussian units in Ref.~\cite[Eq.~(8)]{kim2020quickstart}. The dielectric constant is $\varepsilon_0$ for a relaxed defect, whose screening includes the ions, and $\varepsilon_\infty$ for a vertical transition at frozen nuclei; in silicon the two are equal, 11.7. The bare point-charge term often overcorrects: for the doubly charged carbon vacancy in diamond it does so clearly~\cite[Sect.~III\,B]{komsa2009comparison}, because the defect charge is not a point. The quadrupole term is hard to determine, and a second artifact remains: the average electrostatic potential of a charged cell is set by convention, so that band edges and defect levels of different cells must be aligned. The schemes in current use, by Freysoldt, Neugebauer and Van de Walle and by Kumagai and Oba, model the defect charge as a Gaussian, compute its periodic energy analytically, and align the potential far from the defect, read from the plane-averaged or the atomic-site potentials, with the model potential~\cite[Sect.~II and Eq.~(1)]{kumagai2014electrostatics}. The Kumagai--Oba version also handles an anisotropic dielectric tensor, and shows that the conventional potential-alignment term is already contained in the image-charge correction~\cite[Sect.~IV]{kumagai2014electrostatics}. Without a correction, formation energies of charged defects converge as $1/L_{\mathrm{sc}}$, and their spread between correction schemes translates into defect densities that differ by orders of magnitude~\cite[Sect.~3]{kim2020quickstart}.

\subsection{Color centers}

The optical transition of a color center connects two states of a few localized electrons. Its zero-phonon energy $E_{\mathrm{ZPL}}$ is the difference of the total energies of the relaxed excited and ground states. When both states are well described by single Slater determinants, the excited state is computed by constraining the occupations, promoting an electron from one Kohn--Sham orbital to another and relaxing the structure in that configuration; this is the $\Delta$SCF method. For the negatively charged nitrogen--vacancy center of diamond, with its triplet transition between single-determinant states, a 512-site supercell gives $E_{\mathrm{ZPL}} = \SI{1.757}{eV}$ with a semilocal functional and \SI{2.035}{eV} with the HSE hybrid functional, against the measured \SI{1.945}{eV}~\cite[Sect.~4.1]{alkauskas2014first}. For the T center of silicon, the same approach in a 512-atom supercell with HSE and $\Gamma$-point sampling gives \SI{985}{meV} against \SI{935}{meV}, and the $(0/{-1})$ transition level at \SI{1.07}{eV} above the valence maximum~\cite{dhaliah2022first}. The radiative lifetime, from the transition dipole of the two Kohn--Sham orbitals and the Wigner--Weisskopf rate, comes out as 1.3 to \SI{5}{\micro s}, against a measured \SI{0.9}{\micro s}~\cite{dhaliah2022first}.

The same two relaxed geometries give the vibronic coupling. The Huang--Rhys factor $S_{\mathrm{HR}}$ of \cref{sec:5.6} is the number of phonons emitted on average when the lattice relaxes after the transition, $S_{\mathrm{HR}} = \sum_k\omega_k\Delta Q_k^2/2$, with $\Delta Q_k$ the projection of the mass-weighted displacement between the two minima on the vibrational mode $k$ of frequency $\omega_k$; the zero-phonon fraction of the emission is $\ee^{-S_{\mathrm{HR}}}$. Its convergence with supercell size is slow, because the relaxation reaches into long-wavelength acoustic modes that a small cell does not contain. For the nitrogen--vacancy center, a 64-site supercell gives $S_{\mathrm{HR}} = 3.02$, 20 percent below the converged 3.67, and the converged value needs an embedding scheme: the forces of the transition from a 512-site hybrid-functional calculation, and the vibrational modes of supercells of up to 10\,648 sites from force constants computed with a semilocal functional~\cite[Sects.~2.3 and~4.2]{alkauskas2014first}. The measured value is $3.85 \pm 0.05$~\cite[Sect.~3]{alkauskas2014first}.

$\Delta$SCF has a clear limit. States that are not single determinants, such as the singlet states of the nitrogen--vacancy center, need a correlated treatment of an active space of defect orbitals, embedded in the screened host, and the density transition matrix elements that a dark-matter excitation needs, $|\langle F|n_e(\vect{q})|I\rangle|^2$ over keV momenta, need the many-body states themselves rather than their energies (\cref{sec:16.6}). Shallow donors are the opposite limit: their effective-mass envelope extends over several nanometers, far beyond any supercell, and a supercell calculation can only supply the central-cell correction to the effective-mass theory of \cref{app:D.23}.

\subsection{Displacement thresholds}

The displacement threshold $E_{\mathrm{d}}$ of \cref{sec:7.4} is the smallest kinetic energy that, given to one atom in a given direction, leaves a stable defect. It is computed by molecular dynamics over many directions. The silicon range of 17.5 to \SI{63.5}{eV} is that of the Stillinger--Weber potential, which of three common silicon potentials reproduces the density-functional thresholds best~\cite{kadribasic2018directional}. \Cref{app:E.11} gives the protocol, its cost and its approximations, among them the electronic excitation that Born--Oppenheimer dynamics leaves out (\cref{chap:19}).

\subsection{Tunneling defects in metals}

A defect that tunnels between two configurations $L$ and $R$, such as a light interstitial or an atom in a disordered film, couples to the conduction electrons of a metal through the change of their scattering potential. That coupling damps the tunneling and is described by the ohmic constant $\alpha_{\mathrm{ohm}}$ of \cref{sec:16.8}. In the Born limit, $\alpha_{\mathrm{ohm}} = 4K_{\mathrm{b}}$ with $K_{\mathrm{b}} = (\mathcal{N}_{\mathrm{F}}V_{\mathrm{d}}\Omega)^2$, where $\mathcal{N}_{\mathrm{F}}$ is the density of states per spin and volume at the Fermi level and $V_{\mathrm{d}}$ half the difference of the contact potentials of the two configurations, which sit at $\pm V_{\mathrm{d}}$~\cite[Eq.~(4)]{maillet2023nanomechanical}. Beyond the Born limit the change of potential is described by the scattering matrices $S_L$ and $S_R$ of the conduction electrons at the Fermi energy in the two configurations, with $S = \ee^{2\ii\delta}$ for a channel of phase shift $\delta$, and by the eigenvalues $\ee^{2\ii\tilde\delta_\nu}$ of $S_RS_L^\dagger$, so that $\tilde\delta_\nu$ is the difference of the phase shifts in channel $\nu$. A channel whose phase shift changes by $\tilde\delta$ reduces the long-time overlap of the two Fermi seas as $t^{-(\tilde\delta/\pi)^2}$, the orthogonality exponent of Anderson~\cite{anderson1967infrared}, and an ohmic bath of coupling $\alpha_{\mathrm{ohm}}$ produces $t^{-2\alpha_{\mathrm{ohm}}}$. Equating the two over all channels and both spins,
\begin{equation}
  \alpha_{\mathrm{ohm}} = \frac{1}{2\pi^2}\sum_\sigma\sum_\nu\tilde\delta_\nu^{\,2} ,
  \label{eq:E.10-ohm}
\end{equation}
which in the Born limit, with one channel and phase shifts $\mp\pi\mathcal{N}_{\mathrm{F}}V_{\mathrm{d}}\Omega$ in the two configurations, so that $\tilde\delta = 2\pi\mathcal{N}_{\mathrm{F}}V_{\mathrm{d}}\Omega$, returns $\alpha_{\mathrm{ohm}} = 4(\mathcal{N}_{\mathrm{F}}V_{\mathrm{d}}\Omega)^2 = 4K_{\mathrm{b}}$. A first-principles calculation needs the two relaxed configurations, from a nudged-elastic-band search, identical relaxation of the host around both, and the scattering phases at the Fermi energy, either from a Green's-function method or from the overlap of the two many-electron ground states of a supercell as a function of its size. The orthogonality exponent appears only with fixed occupations: $\vect{k}$ sampling and smearing of the Fermi surface blur the logarithm on which it rests and must be removed or extrapolated.

\begin{workedexample}{A charged defect in a silicon supercell}\label{we:E-sidefect}
We estimate the image-charge error for a singly charged defect, $z_{\mathrm{D}} = \pm1$, in cubic supercells of silicon built from the conventional cell, $L_{\mathrm{sc}} = na$ with $a = \SI{5.431}{\angstrom}$, and with $\alpha\hbar c = \SI{14.40}{eV.\angstrom}$ and $\varepsilon_0 = 11.7$. \Cref{eq:E.10-mp} gives
\begin{equation*}
  E_{\mathrm{corr}} = \frac{2.8373\times\SI{14.40}{eV.\angstrom}}{2\times11.7\times L_{\mathrm{sc}}}
  = \SI{0.161}{eV},\ \SI{0.107}{eV},\ \SI{0.080}{eV}
\end{equation*}
for 64, 216 and 512 atoms ($n = 2$, 3, 4). Even the 512-atom cell, the size used for the T center, carries a point-charge correction of \SI{80}{meV}, larger than the \SI{50}{meV} by which its computed zero-phonon line misses the measurement. Since the point-charge term can overcorrect (above), production calculations use the Gaussian-charge schemes; what the estimate shows is that the error does not become negligible at any affordable size. It falls only as $1/L_{\mathrm{sc}}$, that is, as $N_{\mathrm{at}}^{-1/3}$, so that halving it means eight times more atoms. The cost of a hybrid-functional calculation grows faster than the cube of the number of atoms, and a 4096-atom cell is out of reach where a 512-atom cell is routine; the correction, not the cell, is what converges the result. The neutral T center needs none of this, and its error budget is the functional; its $(0/{-1})$ level at \SI{1.07}{eV} does, since the charged state enters it~\cite{dhaliah2022first}.
\end{workedexample}

